\subsection{Eigenvalues, Eigenspaces, and Spectral Criteria}

Eigenstructure begins by identifying the scalars for which a shifted operator has a nontrivial kernel.

Throughout this subsection, $\mathcal{V}$ is a vector space over $\mathbb{F}$, and $n\in\mathbb{Z}_{>0}$ whenever an $n\times n$ matrix occurs.

\begin{definition} \label{def:eigenvalue-eigenvector} \hlt{Eigenvalue and Eigenvector}\\
Let $T\in\LL(\mathcal{V})$.
\begin{enumerate}[label=(\roman*)]
\setlength{\itemsep}{0pt}
\item A scalar $\lambda\in\mathbb{F}$ is an \hlt{eigenvalue} of $T$ iff $(\exists\,\mathbf{v}\in\mathcal{V}\setminus\{\mathbf{0}\})\,T\mathbf{v}=\lambda\mathbf{v}$.
\item Given $\lambda\in\mathbb{F}$, a vector $\mathbf{v}\in\mathcal{V}$ is an \hlt{eigenvector} corresponding to $\lambda$ iff $\mathbf{v}\neq\mathbf{0}\land T\mathbf{v}=\lambda\mathbf{v}$.
\end{enumerate}
\end{definition}

\begin{definition} \label{def:eigenspace} \hlt{Eigenspace}\\
Let $T\in\LL(\mathcal{V})$ and $\lambda\in\mathbb{F}$. The \hlt{eigenspace} of $T$ corresponding to $\lambda$ is $E_{\lambda}(T)\defeq\ker(T-\lambda I_{\mathcal{V}})=\{\mathbf{v}\in\mathcal{V}:T\mathbf{v}=\lambda\mathbf{v}\}$.\\
Thus $\lambda$ is an eigenvalue of $T$ iff $E_{\lambda}(T)\neq\{\mathbf{0}\}$.
\end{definition}

\begin{definition} \label{def:spectrum-and-geometric-multiplicity} \hlt{Spectrum and Geometric Multiplicity}\\
Let $T\in\LL(\mathcal{V})$.
\begin{enumerate}[label=(\roman*)]
\setlength{\itemsep}{0pt}
\item The \hlt{spectrum} of $T$ is $\operatorname{spec}(T)\defeq\{\lambda\in\mathbb{F}:E_{\lambda}(T)\neq\{\mathbf{0}\}\}$.
\item If $\mathcal{V}$ is finite-dimensional and $\lambda\in\operatorname{spec}(T)$, then the \hlt{geometric multiplicity} of $\lambda$ is $\dim E_{\lambda}(T)$.
\item For $A\in\operatorname{Mat}(n,n,\mathbb{F})$ and $\lambda\in\mathbb{F}$, set $E_{\lambda}(A)\defeq\ker(A-\lambda I_n)$. Define $\operatorname{spec}(A)\defeq\{\mu\in\mathbb{F}:E_{\mu}(A)\neq\{\mathbf{0}\}\}$.
\end{enumerate}
\end{definition}

\begin{theorem} \label{thm:equivalent-eigenvalue-conditions} \hlt{Equivalent Conditions for an Eigenvalue}\\
Let $\mathcal{V}$ be finite-dimensional, let $T\in\LL(\mathcal{V})$, and let $\lambda\in\mathbb{F}$. The following are equivalent:
\begin{enumerate}[label=(\roman*)]
\setlength{\itemsep}{0pt}
\item $\lambda$ is an eigenvalue of $T$;
\item $T-\lambda I_{\mathcal{V}}$ is not injective;
\item $T-\lambda I_{\mathcal{V}}$ is not surjective;
\item $T-\lambda I_{\mathcal{V}}$ is not invertible.
\end{enumerate}
\begin{proof}
\begin{enumerate}[label=(\roman*)]
\setlength{\itemsep}{0pt}
\item[$(i)$] $\Rightarrow (ii)$ By condition (i) and Definition \ref{def:eigenspace}, $\ker(T-\lambda I_{\mathcal{V}})=E_{\lambda}(T)\neq\{\mathbf{0}\}$.\\
The operator $T-\lambda I_{\mathcal{V}}$ is not injective by Proposition \ref{prop:injective-iff-kernel-zero}.
\item[$(ii)$] $\Rightarrow (iii)$ The operator $T-\lambda I_{\mathcal{V}}$ is not surjective by Theorem \ref{thm:finite-dimensional-operator-invertible-equivalences}(ii),(iii).
\item[$(iii)$] $\Rightarrow (iv)$ The operator $T-\lambda I_{\mathcal{V}}$ is not invertible by Theorem \ref{thm:finite-dimensional-operator-invertible-equivalences}(iii),(i).
\item[$(iv)$] $\Rightarrow (i)$ The operator $T-\lambda I_{\mathcal{V}}$ is not injective by Theorem \ref{thm:finite-dimensional-operator-invertible-equivalences}(i),(ii).\\
Hence $E_{\lambda}(T)\neq\{\mathbf{0}\}$ by Proposition \ref{prop:injective-iff-kernel-zero} and Definition \ref{def:eigenspace}, so $\lambda$ is an eigenvalue by Definition \ref{def:eigenvalue-eigenvector}.
\end{enumerate}
\hlt{Structural Commentary}: The eigenspace criterion enters the finite-dimensional injective--surjective--invertible cycle.\\
\hlt{Key Insight}: Spectral values are precisely the singular shifts of $T$.
\end{proof}
\end{theorem}

\begin{theorem} \label{thm:distinct-eigenvalues-independent-eigenvectors} \hlt{Eigenvectors for Distinct Eigenvalues Are Independent}\\
Let $m\in\mathbb{Z}_{>0}$ and $T\in\LL(\mathcal{V})$. Suppose $\lambda_1,\ldots,\lambda_m$ are distinct eigenvalues of $T$, with corresponding eigenvectors $\mathbf{v}_1,\ldots,\mathbf{v}_m$.\\
Then $(\mathbf{v}_1,\ldots,\mathbf{v}_m)$ is linearly independent.
\begin{proof}
Suppose list is dependent. Vector $\mathbf{v}_1$ is nonzero by Definition \ref{def:eigenvalue-eigenvector}(ii), so $\exists$ minimal $k\geq2$ such that $(\mathbf{v}_1,\ldots,\mathbf{v}_k)$ is dependent.\\
Then $(\mathbf{v}_1,\ldots,\mathbf{v}_{k-1})$ is independent and $\mathbf{v}_k\in\operatorname{span}(\mathbf{v}_1,\ldots,\mathbf{v}_{k-1})$. Choose $a_1,\ldots,a_{k-1}\in\mathbb{F}$ satisfying $\mathbf{v}_k=\sum\limits_{i=1}^{k-1}a_i\mathbf{v}_i$.\\
By linearity, eigenvector identities, $\mathbf{0}=(T-\lambda_kI_{\mathcal{V}})\mathbf{v}_k=\sum\limits_{i=1}^{k-1}a_i(\lambda_i-\lambda_k)\mathbf{v}_i$. Thus $(\forall\,i\in\{1,\ldots,k-1\})\,a_i(\lambda_i-\lambda_k)=0$ by independence.\\
The equality $a_1=\cdots=a_{k-1}=0$ holds by distinctness of the eigenvalues, so $\mathbf{v}_k=\mathbf{0}$, contrary to Definition \ref{def:eigenvalue-eigenvector}(ii).\\
\hlt{Structural Commentary}: A minimal dependent prefix and its shifted image isolate every earlier coefficient.\\
\hlt{Key Insight}: Distinct eigenvalue differences cannot hide a nonzero dependence coefficient.
\end{proof}
\end{theorem}

\begin{corollary} \label{cor:number-of-distinct-eigenvalues} \hlt{Number of Distinct Eigenvalues}\\
If $\mathcal{V}$ is finite-dimensional and $T\in\LL(\mathcal{V})$, then $T$ has at most $\dim(\mathcal{V})$ distinct eigenvalues.
\begin{proof}
Set $d\defeq\dim(\mathcal{V})$, and choose a basis $\beta$ of length $d$ by Corollary \ref{cor:finite-dimensional-spaces-and-subspaces-have-bases} and Definition \ref{def:dimension}.\\
If $T$ had $d+1$ distinct eigenvalues, choose corresponding eigenvectors $\mathbf{v}_1,\ldots,\mathbf{v}_{d+1}$. The list $(\mathbf{v}_1,\ldots,\mathbf{v}_{d+1})$ is independent by Theorem \ref{thm:distinct-eigenvalues-independent-eigenvectors}. Since $\beta$ spans $\mathcal{V}$, the inequality $d+1\leq d$ follows by Theorem \ref{thm:finitedimlenofvec}, a contradiction.\\
\hlt{Structural Commentary}: Distinct eigenvalues produce an independent list constrained by a basis-length bound.\\
\hlt{Key Insight}: One independent eigenvector witnesses each distinct eigenvalue.
\end{proof}
\end{corollary}

\begin{corollary} \label{cor:invertible-iff-zero-not-eigenvalue} \hlt{Invertibility and the Zero Eigenvalue}\\
Let $\mathcal{V}$ be finite-dimensional and let $T\in\LL(\mathcal{V})$. Then $T$ is invertible iff $0\notin\operatorname{spec}(T)$.
\begin{proof}
\begin{enumerate}[label=\roman*.]
\setlength{\itemsep}{0pt}
\item[$(\Rightarrow)$] If $T$ is invertible, then $0$ is not an eigenvalue by Theorem \ref{thm:equivalent-eigenvalue-conditions}(iv),(i), so $0\notin\operatorname{spec}(T)$ by Definition \ref{def:spectrum-and-geometric-multiplicity}(i).
\item[$(\Leftarrow)$] If $0\notin\operatorname{spec}(T)$, then $0$ is not an eigenvalue by Definition \ref{def:spectrum-and-geometric-multiplicity}(i), so $T=T-0I_{\mathcal{V}}$ is invertible by Theorem \ref{thm:equivalent-eigenvalue-conditions}(i),(iv).
\end{enumerate}
\hlt{Structural Commentary}: Specializing the shifted-operator criterion at zero identifies invertibility with spectral exclusion.\\
\hlt{Key Insight}: The zero eigenspace is exactly $\ker T$.
\end{proof}
\end{corollary}

\begin{corollary} \label{cor:matrix-invertible-iff-zero-not-spectrum} \hlt{Matrix Invertibility and the Spectrum}\\
Let $A\in\operatorname{Mat}(n,n,\mathbb{F})$. Then $A$ is invertible iff $0\notin\operatorname{spec}(A)$.
\begin{proof}
\begin{enumerate}[label=\roman*.]
\setlength{\itemsep}{0pt}
\item[$(\Rightarrow)$] If $A$ is invertible, then $\operatorname{rank}(A)=n$ by Corollary \ref{cor:square-matrix-invertible-iff-full-rank}, so $\dim(\ker(A))=0$ by Corollary \ref{cor:matrix-rank-nullity}(i).\\
As $\{\mathbf{0}\}\leq\ker(A)$ and both spaces have dimension zero, $\ker(A)=\{\mathbf{0}\}$ by Proposition \ref{prop:dimension-bounds-and-basis-criteria}(ii);  $0\notin\operatorname{spec}(A)$ by Definition \ref{def:spectrum-and-geometric-multiplicity}(iii).
\item[$(\Leftarrow)$] If $0\notin\operatorname{spec}(A)$, then $\ker(A)=\{\mathbf{0}\}$ by Definition \ref{def:spectrum-and-geometric-multiplicity}(iii), so $\dim(\ker(A))=0$ by Definition \ref{def:dimension}.\\
Thus $\operatorname{rank}(A)=n$ by Corollary \ref{cor:matrix-rank-nullity}(i), and $A$ is invertible by Corollary \ref{cor:square-matrix-invertible-iff-full-rank}.
\end{enumerate}
\hlt{Structural Commentary}: Full rank and rank--nullity translate matrix invertibility into zero-kernel spectral language.\\
\hlt{Key Insight}: Zero is spectral exactly when the homogeneous system has a nonzero solution.
\end{proof}
\end{corollary}
